\documentclass[11pt,twocolumn]{article}

\usepackage[margin=0.8in]{geometry}
\usepackage{amsmath,amssymb}
\usepackage{booktabs}
\usepackage{array}
\usepackage{xcolor}
\usepackage{fancyvrb}
\usepackage{enumitem}
\usepackage{titlesec}
\usepackage{balance}
\usepackage{microtype}
\usepackage[hidelinks]{hyperref}

\definecolor{accent}{RGB}{20,70,120}
\definecolor{win}{RGB}{20,110,60}
\definecolor{bad}{RGB}{160,40,40}
\titleformat{\section}{\large\bfseries\color{accent}}{\thesection}{0.5em}{}
\titleformat{\subsection}{\normalsize\bfseries}{\thesubsection}{0.5em}{}
\titlespacing\section{0pt}{8pt}{4pt}
\titlespacing\subsection{0pt}{6pt}{3pt}

\fvset{fontsize=\scriptsize,frame=single,framesep=4pt,xleftmargin=2pt,xrightmargin=2pt}
\setlist{leftmargin=1.2em,topsep=2pt,itemsep=1pt}
\emergencystretch=2em  % absorb tiny overfulls from long \texttt{} identifiers

\newcommand{\best}[1]{\textbf{#1}}
\newcommand{\gn}{{\normalfont\scshape GraphNative}}
\newcommand{\mcp}{{\normalfont\scshape Graph-MCP}}
\newcommand{\plain}{{\normalfont\scshape Plain}}
\newcommand{\hit}[1]{\textcolor{win}{#1}}
\newcommand{\miss}[1]{\textcolor{bad}{#1}}

\title{\vspace{-1.4cm}\bfseries Why Graph-Native Beats Graph-as-a-Tool:\\[2pt]
\large A Case Study in Code-Impact Analysis on Apache Hadoop}
\author{}
\date{}

\begin{document}
\twocolumn[
  \begin{@twocolumnfalse}
  \maketitle
  \vspace{-1.2cm}
  \begin{abstract}
  \noindent
  A coding agent can use a code graph in two ways: as an \emph{optional MCP tool} it chooses to
  call and must rank in-context, or \emph{natively}---with the harness querying and ranking the
  graph before the agent's first token. We show, with real run data on Apache Hadoop
  (14{,}574 Java files), that these are not equivalent. Holding the model
  (\texttt{claude-sonnet-4-6}), the repository, the prompt, and \emph{the underlying graph engine}
  fixed across arms---so the only variable is \emph{where ranking happens}---a graph-native harness
  attains a mean impact-analysis $F_1$ of \best{0.79} versus \best{0.50} for the same agent with the
  same graph attached over MCP and \best{0.56} for a tool-free baseline, at three runs per cell with
  non-overlapping confidence intervals, while being the cheapest arm on every task. We dissect two
  tasks at the level of \emph{which files each arm actually named}. On a hub class
  (\texttt{Resource}, an 982-file blast radius) the MCP arm \emph{collapses} to $F_1\,0.07$: it
  surfaces only 2 of 25 true dependents because it must rank a firehose mid-reasoning, while the
  native harness surfaces 13 of 25 from a pre-ranked draft. On a smaller task
  (\texttt{AbfsOutputStream}) all three arms find every true dependent, but only the native arm
  excludes the test-file noise that drags the others to $F_1\,0.67$---it commits exactly the 10 gold
  files, $F_1\,1.00$. The lesson is concrete: the graph's value is unlocked by \emph{relocating
  ranking out of the model's context and into the harness}, not by exposing a better tool.
  \vspace{0.35cm}
  \end{abstract}
  \end{@twocolumnfalse}
]

\section{The Question, Made Concrete}

Before an agent changes a method it must know what depends on it. A code graph answers this
exactly---\texttt{impact(symbol)} returns the reverse-reachability set. The catch: on a real
codebase that set is a \emph{firehose}. For \texttt{Resource} in Hadoop it is 982 files at
recall~1.0 but precision~0.06. The hard part is not \emph{retrieving} the dependents; it is
\emph{ranking} the 25 that matter above the 957 that don't, under a budget.

This paper asks where that ranking should happen. Two designs:
\begin{itemize}
  \item \mcp{} (graph-as-a-tool): the agent is given the code graph as an MCP server. It decides
  when to call \texttt{impact}/\texttt{callers}, and must rank the returned firehose \emph{in its
  own context, while also reasoning about the task}.
  \item \gn{} (graph-native harness): \emph{before} the agent's first token, the harness runs the
  graph queries, ranks the result with a deterministic structural ranker, drops test files, and
  hands the agent a short, mostly-correct draft to \emph{refine}.
\end{itemize}
Both use the identical graph engine (\textsc{codegraph}; tree-sitter\,$\to$\,SQLite). The only
difference is \emph{when and where} ranking occurs. We show this difference is decisive.

\paragraph{Setup.} All arms run the same model (\texttt{claude-sonnet-4-6}) headlessly through the
same \texttt{claude -p} invocation on a Claude subscription (no API key), differing only in harness
wiring. Tasks are 12 impact questions derived from real merged Hadoop PRs; the gold set $G$ for each
is the actual caller/dependent files touched in that PR. The metric is $F_1$ of the agent's
committed \texttt{dependent\_files} list against $G$ (basename set-match), truncated to a budget of
$k{=}20$---so dumping the firehose tanks precision, and the task becomes \emph{ranking the right few}.
The raw \texttt{impact} query, scored identically, is an oracle \emph{floor} every graph arm must
clear. We report three runs per cell with 95\% confidence intervals.

\section{Headline: Three Runs, Three Tasks}

Table~\ref{tab:main} is the full matrix. \gn{} wins all three tasks on mean $F_1$, with
non-overlapping CIs against the runner-up on the two decisive tasks, \emph{and} is the cheapest arm
everywhere.

\begin{table}[t]\centering\footnotesize
\setlength{\tabcolsep}{3.5pt}
\renewcommand{\arraystretch}{1.12}
\caption{Impact-analysis $F_1$ at budget $k{=}20$, mean\,$\pm$\,95\% CI over $n{=}3$ runs, single
\texttt{claude} gateway. \textsc{Cost} is mean USD/run; \textsc{R/G} is mean Read/Grep calls. Best
per row bold.}\label{tab:main}
\begin{tabular}{lcccc}
\toprule
 & \plain{} & \mcp{} & \best{\gn{}} & oracle \\
\midrule
\multicolumn{5}{l}{\textit{I1 \texttt{Resource} --- hub, 982-file blast radius (gold 25)}}\\
$F_1$    & 0.27\,$\pm$.11 & 0.07\,$\pm$.06 & \best{0.56\,$\pm$.06} & 0.00 \\
cost     & \$0.53 & \$0.38 & \best{\$0.27} & --- \\
R/G      & 0/13 & 0/1 & \best{0/0} & --- \\
\midrule
\multicolumn{5}{l}{\textit{I8 \texttt{AbfsClient} --- clean caller tree (gold 19)}}\\
$F_1$    & 0.74\,$\pm$.29 & 0.77\,$\pm$.00 & \best{0.82\,$\pm$.00} & 0.36 \\
cost     & \$0.57 & \$0.38 & \best{\$0.28} & --- \\
R/G      & 0/28 & 0/0 & \best{0/0} & --- \\
\midrule
\multicolumn{5}{l}{\textit{I10 \texttt{AbfsOutputStream} (gold 10)}}\\
$F_1$    & 0.67\,$\pm$.00 & 0.67\,$\pm$.00 & \best{1.00\,$\pm$.00} & 0.47 \\
cost     & \$0.71 & \$0.39 & \best{\$0.31} & --- \\
R/G      & 0/29 & 0/0 & \best{0/0} & --- \\
\midrule
\best{Mean $F_1$} & 0.56 & 0.50 & \best{0.79} & 0.28 \\
\bottomrule
\end{tabular}
\end{table}

Two facts deserve emphasis before the case studies. First, \gn{} does \best{zero} file reads and
\best{zero} greps on every task: the graph fully displaces filesystem exploration. \plain{} grep-storms
(13--29 greps/task). Second, the arms' \emph{variance} differs: \gn{} and \mcp{} are near-deterministic
($\pm.00$ on I8/I10), while \plain{} is high-variance ($\pm.29$ on I8)---grep-storm success is
luck-dependent. Robustness is not only a higher mean; it is a tighter one.

We now open the two decisive tasks and look at exactly which files each arm named.

\section{Case Study 1 --- The Hub (\texttt{Resource})}

\texttt{Resource} is a YARN god-object: 982 files reference it transitively. The PR
(YARN-11964, clamp negative values) touched 25 of them. This is the pure ranking task---recall is
trivial (the graph contains all 25), precision under budget is everything.

\paragraph{What the arms named (top 20, run r101).}
The gold set (25): {\scriptsize\texttt{Resources, DominantResourceCalculator, AbstractLeafQueue,
ResourceUsage, ResourceCalculator, DefaultResourceCalculator, AbstractCSQueue, QueueMetrics,
FSAppAttempt, FSQueue, AbstractYarnScheduler, RLESparseResourceAllocation, QueueResourceQuotas,
\ldots}}

\begin{Verbatim}[label={Plain (control) -- 74 files; 6/25 gold in top 20}]
LightWeightResource, ResourcePBImpl, *Resources,
*DefaultResourceCalculator, *DominantResource-
Calculator, FifoScheduler, FairScheduler,
FSLeafQueue, FSParentQueue, *FSQueue,
FSQueueMetrics, ConfigurableResource,
*AbstractLeafQueue, CapacitySchedulerConfig,
CSQueueMetrics, ..., *QueueMetrics
   (* = one of the 25 true dependents)
\end{Verbatim}

\begin{Verbatim}[label={Graph-MCP -- 26 files; only 2/25 gold in top 20}]
LightWeightResource, ResourcePBImpl, *Resources,
TestResource, CapacityScheduler, *FSQueue,
DominantResourceFairnessPolicy, AppInfo,
NodeManager, ContainerManagerImpl,
RMContainerAllocator, GetResourceProfile-
ResponsePBImpl, ServiceApiUtil, NodeManager,
TestFifoScheduler, TestFairScheduler, ...
\end{Verbatim}

\begin{Verbatim}[label={GraphNative -- 35 files; 13/25 gold in top 20}]
LightWeightResource, ResourcePBImpl, *Resources,
*DominantResourceCalculator, *DefaultResource-
Calculator, *ResourceCalculator, *ResourceUsage,
*AbstractLeafQueue, *FSAppAttempt, *AbstractCSQueue,
CapacitySchedulerConfig, *QueueMetrics,
*AbstractYarnScheduler, *FSQueue, FairScheduler,
*RLESparseResourceAllocation, RMContainerRequestor,
RMContainerAllocator, *QueueResourceQuotas, RMNodeImpl
\end{Verbatim}

\paragraph{Reading the result.} All three arms see the \emph{same} 982-file graph. What differs is
what reaches the top 20:
\begin{itemize}
  \item \plain{} (\textcolor{bad}{$F_1$ 0.27}) brute-forces with 13 greps, names \emph{74} files, and
  buries the 6 true dependents it does find among scheduler/config noise. Over-enumeration is a
  precision tax.
  \item \mcp{} (\textcolor{bad}{$F_1$ 0.07}) is the cautionary case. It \emph{has} the graph and
  calls it (7 graph calls), but having pulled a 982-file firehose into context it cannot rank it
  while also reasoning, so it commits a short, hedged list seeded by surface name-matches and
  unrelated managers (\texttt{NodeManager}, \texttt{AppInfo}, \texttt{TestResource})---landing only
  \emph{2 of 25}. The tool was present; the ranking was absent.
  \item \gn{} (\textcolor{win}{$F_1$ 0.56}) starts from the harness's pre-ranked, test-free draft and
  refines it. Its top 20 contains \emph{13 of 25} true dependents---the dense, genuinely-coupled
  resource/queue/scheduler classes---because reference-density and caller ranking were computed in
  code \emph{before} the agent had to think.
\end{itemize}
Same model, same graph, same budget. The $8\times$ gap between \gn{} and \mcp{} is entirely
\emph{where the firehose got ranked}. This is the paper's thesis in one task.

\section{Case Study 2 --- The Test-File Trap (\texttt{AbfsOutputStream})}

The second task is the opposite regime: a \emph{small} blast radius (gold 10). Here recall is easy
for everyone---and indeed all three arms find all 10 true dependents. The difference is what
\emph{else} they put in the top 20.

\begin{table}[h]\centering\footnotesize
\setlength{\tabcolsep}{5pt}
\renewcommand{\arraystretch}{1.1}
\caption{I10: every arm achieves recall 1.0 (all 10 gold found). Precision is decided by false
positives in the top 20.}\label{tab:i10}
\begin{tabular}{lccc}
\toprule
 & gold found & false-pos & $F_1$ \\
\midrule
\plain{}        & 10/10 & 10 (5 tests) & 0.67 \\
\mcp{}          & 10/10 & 10 (7 tests) & 0.67 \\
\best{\gn{}}    & \best{10/10} & \best{0} & \best{1.00} \\
\bottomrule
\end{tabular}
\end{table}

The false positives are revealing. \plain{} pads with \texttt{TestAbfsOutputStream},
\texttt{ITestAbfsOutputStream}, \texttt{AbstractAbfsIntegrationTest}, \ldots; \mcp{} pads with
\texttt{ITestAzureBlobFileSystemChecksum}, \texttt{ITestAbfsOutputStreamStatistics}, \ldots\ Both
arms let test files---which exhibit the \emph{highest} reference density on a real codebase, and are
\emph{never} the answer in a production fix---consume budget slots, halving precision to 0.50.

\gn{} commits \best{exactly the 10 gold files, zero false positives}. The reason is a single
deterministic rule in the harness ranker: \emph{test-file demotion}. Files matching
\texttt{Test*}/\texttt{*Test}/\texttt{ITest*}/\texttt{/test/} are pushed below every production
candidate before the draft is built. The agent never has to relearn, mid-task, that
\texttt{TestAbfsOutputStream} is not a dependent---the harness already knew.

\section{Why the Ranker, Not the Tool, Is the Win}

The two case studies isolate the mechanism. We can also measure it directly, with no agent in the
loop, by scoring the raw graph output against gold (Table~\ref{tab:ranker}). The raw \texttt{impact}
firehose scores $F_1\,0.169$ (held-out, $n{=}9$). A purely structural ranker over the \emph{same}
graph---reference density, an additive direct-caller bonus, type-name inheritance match, and
test-file demotion---triples it to \best{0.519}, beating the raw-graph floor on all 9 held-out tasks.
This ranker is what the harness injects at turn~0; it is gold-blind and validated on tasks never used
for tuning.

\begin{table}[h]\centering\footnotesize
\setlength{\tabcolsep}{6pt}
\renewcommand{\arraystretch}{1.1}
\caption{Held-out ranking quality ($n{=}9$, never tuned), no agent. The ranker---a harness
component---is where the graph's value is realized.}\label{tab:ranker}
\begin{tabular}{lc}
\toprule
 & mean $F_1$ \\
\midrule
raw \texttt{impact} firehose (floor) & 0.169 \\
\;+ structural ranker (turn-0) & \best{0.519} \\
\bottomrule
\end{tabular}
\end{table}

The point is that \mcp{} \emph{has access to the identical \texttt{impact} query}---it simply has to
rank the output in-context, which Case Study~1 shows it does poorly on a hub. Moving that ranking
into deterministic harness code, run before the agent's first token, is the entire delta. It is a
\emph{harness} property, not a \emph{tool} property.

\section{An Honest Foil: potpie}

To check that this is a property of native integration rather than of \textsc{codegraph}
specifically, we ran the same gold against a second graph product, \emph{potpie} (a Neo4j +
tree-sitter codebase-intelligence system). Two findings, both disclosed plainly:
\begin{itemize}
  \item \textbf{Operability.} potpie's LLM agents require a paid API key and a multi-service Docker
  stack; they cannot run on a Claude subscription. Its \emph{graph} can be built LLM-free, which is
  what we benchmarked. \textsc{codegraph}'s whole graph-native arm, by contrast, runs on the
  subscription from a single local binary.
  \item \textbf{Graph quality vs.\ ranking.} potpie's raw name-reference blast radius actually
  \emph{out-recalls} \textsc{codegraph}'s raw \texttt{impact} at top-20 (held-out $F_1$ 0.383 vs
  0.169)---a creditable graph. But \textsc{codegraph}'s \emph{ranked} output (0.519) beats both raw
  graphs, and on the \texttt{Resource} hub potpie's 603-file \emph{unranked} set scores $F_1\,0.00$
  at top-20 while \gn{} scores 0.56. (Caveat: potpie's set is undirected name-reference---higher
  recall, lower precision than a directed query---and unranked, so this is a statement about raw
  recall, not delivered quality.)
\end{itemize}
The cross-product comparison reinforces the thesis: a good raw graph is necessary but not
sufficient; the \emph{ranking and pre-computation layer}---a harness responsibility---is what turns a
graph into answers.

\section{Threats to Validity}
\textbf{Construct.} These tasks measure impact \emph{retrieval} (a ranked file list), not code
synthesis; the draft-to-refine mechanism is partly self-fulfilling within this construct, and
transferring the result to a diff-producing, test-pass-scored task is the natural next step.
\textbf{External.} One Java repository; the structural signals reflect general dependency-graph and
JVM naming conventions but cross-language generalization is unmeasured. \textbf{Scope.} $n{=}3$ over
three tasks spanning the difficulty range; the aggregate gap is more robust than any single cell, and
no arm sweeps---\plain{} and \mcp{} are within noise of each other and of \gn{} on the clean task,
which is the honest, discriminative picture.

\section{Conclusion}
Exposing a code graph to an agent as an MCP tool under-uses it: the agent must both choose to call it
and rank a firehose mid-reasoning, and on a hub class it ranks that firehose so poorly it scores
$8\times$ \emph{below} a native harness using the same graph. Treating the graph as the harness's
primary retrieval surface---queried and ranked before the first token, test files removed, handed
over as a draft to refine---wins all three tasks at lower cost and tighter variance, on a held-out,
gaming-resistant, multi-run benchmark. The gain comes not from a bigger model or a better tool, but
from relocating retrieval and ranking out of the model's context and into deterministic code where
the structure of the dependency graph can be exploited directly.

\balance
\end{document}
